\documentclass{alex_summary}

\begin{document}
\section*{Reading summary}

In \blockcquote{cirac_goals_2012}{Goals and opportunities in quantum simulation}, a \textbf{brief overview} of the research field of \textbf{quantum simulation} is given, outlining basic concepts, possible application, and challenges. The paper \blockcquote{houck_-chip_2012}{On-chip quantum simulation with superconducting circuits} introduces the reader to quantum simulations using the \textbf{superconducting circuits platform}. The authors then focus on possible applications ranging from condensed matter physics to quantum chemistry and extend the discussion to include experimental implementations.

The idea behind a quantum simulator, as opposed to a quantum computer, is to \textbf{mimic} the behaviour of a \textbf{quantum system} such that quantities of interest are identical between the simulator and the system. The key concept here is the \textbf{matching of Hamiltonians} between both systems. By being able to switch or \textbf{tune interactions}, one can study the effect of sweeping parameters, essentially \textbf{driving phase transitions} in the system. Another possibility enabled by quantum simulators is referred to as \textbf{quantum annealing}, where one investigates the properties of ground states by \textbf{adiabatically transforming} one Hamiltonian into another.

Simulating the behaviour of quantum systems on \textbf{superconducting circuits} has one \textbf{major advantage} over alternatives such as ultracold atoms in optical lattices or trapped ions: unlike atoms, which must obey conservation laws, circuit excitations used as the basis for simulation are \textbf{quasiparticles}, for which particle number need not be conserved (hence the term \blockquote{artificial atoms}). This fact allows superconducting circuits to imitate the behaviour of open systems that couple to a bath with a large number of internal degrees of freedom. One prominent and well-known example is the \textbf{Bose-Hubbard model}. This describes a spinless bosonic system in a lattice with a chemical potential on each lattice site and nearest-neighbour hopping between adjacent sites. Depending on the relative coupling strength between on-site and neighbour interactions, either a \textbf{superfluid or Mott-insulator} state arises. Both cases can be retrieved using theoretical tools (such as mean-field or perturbative approaches), but the \textbf{phase transition} diverges in all treatments. This motivates the use of a quantum simulator in which the coupling terms in the Hamiltonian can be tuned to drive a phase transition.\\

\begin{question1}
	Why is decoherence not as big of a problem in quantum simulators as opposed to quantum computers? On the one hand I get the point that small defects do not significantly affect large systems, on the other hand I know of the butterfly effect which states the exact opposite. Does the first statement maybe only apply for certain types of systems?
\end{question1}

\begin{question3}
	I am a bit confused what role time plays in a quantum simulation. In gate-based computations time is only considered in terms of coherence times but all gate operations have no concept of time. In quantum simulation one influences the system to study some dynamics. One probably has to consider the response time of the simulation platform but does one have to take into account the response time of the system of interest and match both in some way? I dont know if this question makes sense, my line of thought is inspired by Reynolds-number matching in fluid dynamics where one can compensate for different viscosity (e.g. replacing air with water) by changing the fluid velocity.
\end{question3}

\printbibliography
\end{document}
